\chapter{Texto e estrutura}
\label{chap:texto-estrutura}

\section{Parágrafos e ênfase}

No código-fonte, uma quebra de linha comum equivale a um espaço. Uma linha vazia
inicia outro parágrafo. Deixe o LaTeX decidir onde quebrar as linhas; isso mantém
o texto adaptável quando margens ou fontes mudarem.

Use comandos semânticos para destacar conteúdo:

\begin{lstlisting}[language=LaTeXBook]
Este termo e \emph{importante}.
Este trecho esta em \textbf{negrito}.
O comando \LaTeX produz o logotipo do sistema.
\end{lstlisting}

Evite usar negrito e tamanho manual para simular títulos. Os comandos de seção
registram a hierarquia, controlam a numeração e alimentam o sumário automático.

\section{Capítulos, seções e sumário}

A classe \arquivo{book} oferece capítulos. Uma estrutura inicial pode ser:

\begin{lstlisting}[language=LaTeXBook]
\documentclass{book}
\begin{document}
\tableofcontents
\chapter{Introducao}
\section{O problema}
\subsection{Contexto historico}
\chapter{Desenvolvimento}
\end{document}
\end{lstlisting}

O sumário pode aparecer vazio na primeira compilação. LaTeX grava os títulos em
um arquivo auxiliar e os lê na passagem seguinte. \arquivo{latexmk} cuida disso
automaticamente.

Use a hierarquia com intenção. Um capítulo trata de um grande assunto; uma seção
organiza uma etapa desse assunto; uma subseção detalha um aspecto. Muitos níveis
produzem um texto fragmentado e difícil de navegar.

\section{Listas e notas}

Listas não numeradas usam \arquivo{itemize}; sequências usam
\arquivo{enumerate}:

\begin{lstlisting}[language=LaTeXBook]
\begin{itemize}
  \item Planejar o conteudo.
  \item Escrever o primeiro rascunho.
  \item Revisar e compilar.
\end{itemize}

\begin{enumerate}
  \item Criar a pasta.
  \item Criar o arquivo principal.
  \item Gerar o PDF.
\end{enumerate}
\end{lstlisting}

Notas de rodapé são inseridas no próprio ponto da chamada com
\comando{footnote\{Texto da nota.\}}. Use-as para observações complementares,
não para esconder explicações essenciais.

\section{Rótulos e referências cruzadas}

Nunca digite manualmente ``veja o capítulo 3''. A numeração pode mudar. Coloque
um rótulo após o elemento e use \comando{ref}:

\begin{lstlisting}[language=LaTeXBook]
\chapter{Metodos}
\label{chap:metodos}

Como veremos no Capitulo~\ref{chap:metodos}, ...
\end{lstlisting}

O caractere \verb|~| cria um espaço que não permite quebra entre a palavra e o
número. Prefixos como \arquivo{chap:}, \arquivo{sec:}, \arquivo{fig:} e
\arquivo{tab:} ajudam a reconhecer cada tipo de rótulo.

\begin{pratica}
  Crie dois capítulos com duas seções cada. Insira um sumário e faça o segundo
  capítulo referenciar uma seção do primeiro. Depois, inverta a ordem dos
  capítulos e confirme que a referência se atualiza.
\end{pratica}

